% !TEX root = ../main.tex
% !TEX program = lualatex
\documentclass[../main.tex]{subfiles}
\IfSubfilesClassLoaded{\externaldocument{\subfix{../build/main}}}{}
% ====================
% File: 18B_test-eval.tex
% ====================
\begin{document}
\IfSubfilesClassLoaded{\chapter{EDO Study Guide}}{}
% === Test & Eval (EDO 3.5.5 and 3.6.6) ===
\section{Test \& Evaluation}
% # 18B_test-eval integration directions.
% 1. Please ensure the following are addressed:
%   - Ensure that all Topic Learning Objectives are addressed in the Study Guide
%   - Define T&E. Difference between a test and an evauluation?
%   - Relationship to verification and validation
%   - What are the statuartory requirements? Which key roles did they make and what do they do?
%   - Describe and define the DoW T&E Organization
%   - Detailed description of Developmental T&Es. Define each one.
%   - Detailed description of Live Fire T&E
%   - Detailed description of Operational T&E (Ininclude EOA, OA, IOT&RE, and FOT&E)
%   - Define the T&E terms
%   - Describe CTPs and how do they relate to TPMs, ICD, CDD, APB, etc.
%   - Detailed description of the \ac{temp}
%   - Describe Modeling and Simulation. When and why is it used?
% 2. Include placeholders for figures at appropriate spots: (1) Test and Evaluation graphic for MCA; (2) DoW T&E Org
% 3. Update `acronyms.def` and `edo.bib` with new entries as needed. Verify no duplicate acronym labels
% 4. Coursebook is the first source document, but relevant sources from the Coursebook, DAU, and other resources should also be used.
% 5. Cross-verify that the information from this section flows logically with the rest of the study guide.
% 6. Review entire study guide and update the Key Roles Appendix.
% 7. Review entire study guide and verify proper usage of acronym calls inside body text (no calls within headings, sections, tables, etc.)
% 8. Ensure consistency with remainder of the Study guide.

% * Enter content here:
\subsection{Board Prep Summary}
\ac{te} provides the objective evidence that a system can achieve the warfighter outcomes promised in requirements documents before obligating the fleet to fielding, sustainment, and congressional reporting commitments~\autocite{edo-3-5-5-test-eval-part-i-2025,DoDI5000-85}. A test is the controlled activity (lab, range, digital thread, fleet event) that produces data, while an evaluation is the analytic judgment that compares that data to thresholds, suitability expectations, and risk acceptance criteria~\autocite{edo-3-5-5-test-eval-part-i-2025}. Effective \ac{te} integrates developmental, operational, and live-fire evidence so that each \ac{mca} milestone is backed by verified design data and validated mission performance insights rather than optimistic modeling alone~\autocite{DoDI5000-89}. Programs that synchronize \ac{te} with systems engineering and program control expose performance gaps early, recycle fixes through the technical baselines, and enter \ac{otrr} with a defendable \ac{temp}, resource plan, and deficiency scrub.

\begin{figure}
  \centering
  \includegraphics[width=0.8\linewidth]{TE_MCA.png}
  \caption[MCA TE Touchpoints]{Key \ac{mca} test and evaluation touchpoints that frame \ac{dt}, \ac{optest}, and certification decisions. \srcCite{edo-3-5-5-test-eval-part-i-2025}.}\label{fig:TE_MCA}
\end{figure}

\subsection{Definition and Relationship to Verification and Validation}
\ac{dte} verifies ``did we build the system right'' by showing the technical baseline meets allocated requirements, standards, and interface controls, while \ac{ote} validates ``did we build the right capability'' by proving Sailors and Marines can accomplish missions under representative threat, environment, and sustainment conditions~\autocite{DoDI5000-02}. Verification relies on repeatable measurements (e.g., propulsion plant efficiency in a land-based engineering site) and feeds configuration control; validation uses operational scenarios, tactics, and logistics supportability to demonstrate mission effectiveness and suitability. \ac{dodi}~5000.88 further directs programs to execute frequent, iterative end-user validation of features and usability so mission thread owners can accept the delivered capability before major decisions~\autocite{DoDI5000-88}. \ac{vv} and \ac{vva} activities must be planned in the \ac{temp} and \ac{sep} so that models, simulations, hardware-in-the-loop facilities, and live events share common data assumptions and confidence levels before milestone reviews~\autocite{DoDI5000-89}.
\textit{Verification} therefore proves compliance with the documented specification (analysis, inspection, test, demonstration), whereas \textit{validation} proves the warfighter accepts the delivered capability and can close the mission threads; both sides of the \ac{vv} matrix must be satisfied before the \ac{mdauthority} can certify readiness for production or fleet release~\autocite{DoDI5000-89}.

\subsection{Operational v.s.\ Developmental Testing}
\begin{description}[style=nextline,labelsep=0.5em,font=\bfseries]
  \item[\textbf{Purpose.}] \ac{dte} is an engineering feedback loop: find design, interface, reliability, cybersecurity, safety, and manufacturability problems early enough to fix the baseline. \ac{ote} is a readiness judgment: determine whether a production-representative system is operationally effective, suitable, survivable, and supportable for Fleet use~\autocite{DoDI5000-89,edo-3-5-6-test-eval-part-ii-2025}.
  \item[\textbf{Primary question.}] Developmental testers ask whether the program built the system correctly against allocated technical requirements. Operational testers ask whether operators can accomplish representative missions with the system under realistic threat, environment, training, logistics, and command-and-control conditions~\autocite{DoDI5000-02}.
  \item[\textbf{Timing and environment.}] \ac{dte} starts early and repeats through component, subsystem, integration, qualification, and certification events in labs, digital environments, ranges, and prototypes. \ac{ote} occurs when the configuration, support package, operators, training, and threat-representative scenarios are mature enough to support an operational adequacy decision before \ac{lrip}, \ac{frp}, Fleet release, or major follow-on fielding~\autocite{DoDI5000-85,USC-10-2399}.
  \item[\textbf{Who leads.}] Program managers, the \ac{cdt}, \ac{ldto}, warfare centers, contractors, and certification authorities normally lead \ac{dte}. Independent Operational Test Agencies such as \ac{comoptevfor}, with \ac{dote} oversight for covered programs, lead \ac{ote} so the Fleet and Congress receive an evaluation independent of program advocacy~\autocite{DoDI5000-89,USC-10-139}.
  \item[\textbf{Evidence and failure meaning.}] \ac{dte} failures are expected risk-discovery events that drive redesign, deficiency correction, reliability growth, or requirements trades. \ac{ote} failures indicate the system may not be ready for the next operational decision; they can delay fielding, restrict Fleet release, require corrective action, or prevent a beyond-\ac{lrip}/\ac{frp} decision~\autocite{USC-10-2399,edo-3-5-6-test-eval-part-ii-2025}.
\end{description}
\note{Board answer: \ac{dte} builds technical confidence while \ac{ote} proves operational fitness. Combined testing may collect shared data, but it does not erase the independence, realism, or adequacy expectations of operational evaluation~\autocite{DoDI5000-02,DoDI5000-89}.}

\subsection{Statutory Requirements and Key Roles}
\begin{description}
  \item[\textbf{10~U.S.C.~\S~139.}] Establishes the \ac{dote} as an independent advisor to \ac{secdef} and Congress; \ac{dote} approves operational test plans, monitors execution, and must issue a public report for covered programs~\autocite{USC-10-139}.
  \item[\textbf{10~U.S.C.~\S~2399.}] Prohibits a program from proceeding beyond \ac{lrip} or requesting Full-Rate Production approval until \ac{dote} and the lead Operational Test Agency submit their independent reports to Congress~\autocite{USC-10-2399}.
  \item[\textbf{10~U.S.C.~\S~2366.}] Requires \ac{lfte} or an approved waiver for covered systems (personnel, platform, and munitions lethality/survivability) before beyond-\ac{lrip} decisions; it also directs reporting of vulnerability/lethality findings to defense committees~\autocite{USC-10-2366}.
  \item[\textbf{\ac{dasdte}.}] Serves as the \ac{osd} focal point for test resources, designates \acp{cdt} for \ac{acat}~I programs, charters \acp{ldto}, and co-chairs Defense-level \ac{te} Working Integrated Product Teams~\autocite{DoDI5000-89}.
  \item[\textbf{\ac{cdt} and \ac{ldto}.}] The \ac{cdt} orchestrates the integrated test schedule, data requirements, and deficiency resolution cadence, while the \ac{ldto} (often a warfare center or system command test directorate) furnishes instrumentation, modeling, and range execution authority for \ac{dte}~\autocite{DoDI5000-89}.
  \item[\textbf{Operational Test Agencies.}] \ac{comoptevfor}, \ac{afotec}, \ac{atec}, and \ac{mcotea} plan and execute \ac{ote} \ac{mda} record~\autocite{edo-3-5-6-test-eval-part-ii-2025}.
\end{description}

\begin{figure}[htbp]
  \centering
  \begin{tikzpicture}[
      node distance=1.5cm and 2.2cm,
      org/.style={draw=Blue, fill=Blue!10, rounded corners=2pt, align=center, inner sep=6pt, minimum width=3.8cm, text width=3.8cm},
      arrow/.style={-Latex, line width=1pt, draw=SteelGray}
    ]
    \node[org] (dote) {OSD DOT\&E\\Statutory adequacy oversight};
    \node[org, right=of dote] (dasdte) {DASW(T\&E)\\Policy, resources, test infrastructure};
    \node[org, below left=0.4cm and 1.4cm of dote] (otas) {Service OTAs\\COMOPTEVFOR / AFOTEC / ATEC / MCOTEA};
    \node[org, below right=0.4cm and 1.4cm of dote] (exec) {Service T\&E executors\\Warfare centers, ranges, labs, M\&S};
    \node[org, below=of otas] (programs) {Component acquisition chain\\PEO / PM / Chief Dev Tester / LDTO};
    \node[org, below=of exec] (fleet) {Warfighter users\\Type commanders / Fleet sponsors};

    \draw[arrow] (dote) -- node[above left]{OT adequacy reports} (otas);
    \draw[arrow] (dasdte) -- node[above right]{T\&E policy / resources} (exec);
    \draw[arrow] (dote) -- node[above]{Congress reports} (programs);
    \draw[arrow] (dasdte) -- node[above]{CDT / LDTO charters} (programs);
    \draw[arrow] (otas) -- node[left]{OT results / suitability} (programs);
    \draw[arrow] (exec) -- node[right]{DT data / certifications} (programs);
    \draw[arrow] (programs) -- node[left]{fielding recommendations} (fleet);
    \draw[arrow] (fleet) -- node[right]{user feedback / mission threads} (programs);
  \end{tikzpicture}
  \caption[DoD TE organization]{\ac{dod} test and evaluation governance chain linking \ac{osd} oversight, Service Operational Test Agencies, warfare centers/ranges, and program-level test leadership.}
  \label{fig:dow-te-org}
\end{figure}

\subsection{Developmental Test Activities}
\begin{description}
  \item[\textbf{Component and qualification testing.}] Bench, lab, and land-based engineering site events that verify materials, software builds, cyber hardening, interoperability, and safety factors before integration; data feeds configuration audits and reliability growth curves~\autocite{edo-3-5-5-test-eval-part-i-2025}.
  \item[\textbf{Subsystem and integration testing.}] Hardware-in-the-loop, systems integration labs, and shore-based platform integration tests that expose interface defects and mission thread timing issues before at-sea trials~\autocite{edo-3-5-5-test-eval-part-i-2025}.
  \item[\textbf{System-level \ac{dte}.}] Full-up prototypes or first-of-class hulls execute combined developmental/operational scenarios, cybersecurity assessments, and logistics demonstrations to validate readiness for \ac{trr} and \ac{techeval}~\autocite{DoDI5000-85}.
  \item[\textbf{Reliability growth and maintainability demonstrations.}] Planned stress events and failure-mode tracking prove that \ac{mtbomf} and corrective maintenance timelines align with sustainment strategies before Operational Test Readiness~\autocite{DoDI5000-89}.
  \item[\textbf{Certification testing.}] Airworthiness, safety, information assurance, weapon certification, and electromagnetic environmental effects tests that deliver statutory certifications needed for fielding~\autocite{edo-3-5-5-test-eval-part-i-2025}.
\end{description}
\note{Edge case: Software pathway efforts may collapse \ac{techeval}/\ac{opeval} into continuous \ac{devsecops} assessments, but the program must still show how combined testing and user evaluations satisfy Title~10 adequacy requirements before fielding~\autocite{DoDI5000-02}.}

\subsection{Live Fire Test and Evaluation}
\begin{description}[style=nextline,labelsep=0.5em,font=\bfseries]
  \item[\textbf{When required.}] Title~10 demands \ac{lfte} for every new or significantly upgraded crew-occupied platform, combat vehicle, major munitions program, or missile defense component that will face threat lethality; full-up system-level shots must be realistically stressed prior to requesting Full-Rate Production approval~\autocite{USC-10-2366,DoDI5000-89}.
  \item[\textbf{When waived.}] Only the \ac{secdef} can waive full-up \ac{lfte} when the system is unmanned, when destructive testing would be unreasonably expensive or impractical, or when alternative testing (component shots, validated modeling, allied data) provides equivalent insights; Congress must be notified of the rationale and the alternative plan~\autocite{USC-10-2366,DoDI5000-89}.
  \item[\textbf{Execution and use.}] Programs plan vulnerability and lethality events in the \ac{temp}, employ full-up articles plus critical subsystems and surrogates, and collect battle damage, crew survivability, and cascading-effect data to inform design hardening, tactics, and survivability equipment upgrades; the results become an annex to the \ac{dote} report before the \ac{frpdr}~\autocite{DoDI5000-89,DoDI5000-85}.
\end{description}

\subsection{Operational Test Phases}
\begin{description}
  \item[\textbf{\ac{eoa}.}] A qualitative review, often using simulations or limited field events, that informs \ac{cdd} refinement and acquisition strategy decisions before Milestone~B~\autocite{edo-3-5-6-test-eval-part-ii-2025}.
  \item[\textbf{\ac{oa}.}] Focused operational events (sometimes embedded in Fleet Battle Problems or Fleet Exercises) that characterize mission effectiveness, training load, and logistics demands prior to \ac{lrip} or dedicated \ac{iote}~\autocite{edo-3-5-6-test-eval-part-ii-2025}.
  \item[\textbf{\ac{iotre}.}] Navy programs often combine \ac{iote} and Readiness Evaluation to prove both mission performance and Fleet introduction readiness; Title~10 requires that the \ac{iote} portion be adequate, realistic, and reported to Congress before \ac{frp}~\autocite{USC-10-2399}.
  \item[\textbf{\ac{fote}.}] Conducted after \ac{frp} or major configuration changes to verify corrective actions, \ac{cpd} thresholds, cybersecurity resilience, and software drops prior to wide deployment~\autocite{DoDI5000-89}.
\end{description}

\subsection{Fundamental Test and Evaluation Terms}
\begin{description}
  \item[\textbf{\ac{coi}.}] The mission-outcome questions (``Can the ship defend itself against a raid?'') that structure \ac{ote} plans and reports~\autocite{edo-3-5-6-test-eval-part-ii-2025}.
  \item[\textbf{\ac{moe} \& \ac{mos}.}] Measures of Effectiveness and Suitability quantify warfighting value (probability of raid annihilation) and fleet readiness attributes (reliability, maintainability, training burden) tied to \acp{coi}~\autocite{edo-3-5-6-test-eval-part-ii-2025}.
  \item[\textbf{\ac{mop}.}] Measures of Performance (e.g., radar detection range) show lower-level technical behaviors that roll up to \acp{moe}/\ac{mos}~\autocite{edo-3-5-5-test-eval-part-i-2025}.
  \item[\textbf{Deficiency categories.}] Category~I deficiencies are mission-critical failures that could cause death, severe injury, or major mission loss; Category~II items are significant but less critical issues requiring tracking and correction~\autocite{edo-3-5-6-test-eval-part-ii-2025}.
  \item[\textbf{Combined testing.}] Planned events that collect developmental and operational data simultaneously to save schedule and align datasets for evaluation teams~\autocite{DoDI5000-02}.
\end{description}

\subsection{Critical Technical Parameters and Technical Performance Measures}
\acp{ctp} are the small subset of design attributes that must be controlled to assure a \ac{kpp} or \ac{ksa} remains achievable; they are derived from the \ac{icd}/\ac{cdd} and baselined inside the \ac{apb}~\autocite{DoDI5000-85}. Each \ac{ctp} maps to one or more \acp{tpm} that engineers track continuously using lab data, digital models, or \ac{dte} results; \acp{tpm} trends highlight when technical debt threatens \ac{temp} and \ac{sep} assumptions so leadership can adjust scope, schedule, or funding. Programs are expected to show a clean trace from \ac{coi}~$\rightarrow$ \ac{kpp}/\ac{ksa}~$\rightarrow$ \ac{ctp}~$\rightarrow$ \ac{tpm}~$\rightarrow$ planned test points in the \ac{temp}, ensuring every test objective underwrites an operational outcome and every requirement has an observable verification method~\autocite{DoDI5000-89,edo-3-5-5-test-eval-part-i-2025}.

\subsection{Test and Evaluation Master Plan}
The \ac{temp} captures the integrated \ac{te} strategy, schedule, resources, and data management approach across the lifecycle; it harmonizes the \ac{sep}, \ac{lcsp}, and \ac{apb} so that every milestone package shows how risk will be retired~\autocite{DoDAcqGuidebook}. Part~I summarizes program context and key performance requirements, Part~II defines the integrated test program (\ac{dt}, \ac{optest}, certification, modeling, and analysis), Part~III outlines resource needs (ranges, threat surrogates, workforce, instrumentation), and annexes address data rights, stats, and cyber accreditation~\autocite{DoDI5000-89}. \ac{acat}~I(D) \acp{temp} require \ac{dote} approval; other \acp{temp} are approved by the Component Acquisition Executive. Updates are mandatory when requirements change, critical deficiencies emerge, or significant schedule/funding revisions alter test sequencing~\autocite{DoDI5000-85}.

\subsection{Modeling and Simulation}
Modeling and simulation provide the only practical way to explore the mission design space before hardware exists, rehearse high-risk scenarios (e.g., hypersonic raid defense), and extend sparse live-test data across the operational envelope~\autocite{edo-3-5-5-test-eval-part-i-2025}. Programs must document model purpose, pedigree, uncertainty, and \ac{vva} status in the \ac{temp}; accredited models can bound live-test needs, size instrumentation, and support decision-quality evaluations when safety, cost, or treaties limit destructive testing~\autocite{DoDI5000-89}. Mission engineering digital threads (system-of-systems simulations, campaign models) allow combined \ac{dt}/\ac{optest} teams to choreograph events, align data management plans, and feed statistical confidence calculations before and after live events.

\subsection{MCA Testing Lifecycle}
\ac{mca} programs must align \ac{te} activities with milestone decision points to provide defensible evidence for phase transitions~\autocite{DoDI5000-85}. DoDI~5000.98, effective December~9,~2024, consolidates \ac{ote} and \ac{lfte} policy previously in DoDI~5000.89, while DoDI~5000.89 remains the primary source for integrated \ac{te} and \ac{dte} policy~\autocite{DoDI5000-98,DoDI5000-89}. Cyber \ac{dte} policy is specified in DoDM~5000.103, effective February~25,~2026~\autocite{DoDM5000-103}.

\begin{longtblr}[
    label   = {tab:mca-te-lifecycle},
    caption = {MCA Testing Lifecycle by Phase},
    entry   = {MCA Testing Lifecycle by Phase},
    remark{\tSource} = {\tabCite{DoDI5000-85,DoDI5000-98,DoDI5000-89,DoDM5000-103}}
  ]
  {colspec = {@{} X[0.85] X[1.6] X[1.4] X[1.1] X[1.3] @{}}}
  \toprule
  {Phase / Decision} & {Primary Test Emphasis} & {Evidence Needed} & {Decision Supported} & {Key Players} \\
  \midrule
  Pre-Milestone~A / \ac{msa} & Technology readiness, feasibility testing, early \ac{eoa} & \ac{trl}~6 evidence, \ac{aoa} data, initial \ac{temp} outline & \ac{mdd} approval, \ac{aoa} guidance & \ac{pm}, \ac{cdt}, warfare centers, \ac{ldto} \\
  \midrule
  \ac{tmrr} & Component and subsystem \ac{dte}, competitive prototyping, cyber \ac{dte} planning & \ac{pdr} completion, technology demonstration, \ac{cddv} inputs, \ac{temp} Part~I & Milestone~A / \ac{drfprd} & \ac{pm}, \ac{cdt}, \ac{ldto}, contractor, \ac{dasdte} \\
  \midrule
  Milestone~B / \ac{emd} entry & System-level \ac{dte}, integration testing, early \ac{oa}, cyber \ac{dte} execution & \ac{cdr} completion, \ac{temp} approval, \ac{apb} baselines, \ac{otrr} readiness & \ac{emd} authorization & \ac{mdauthority}, \ac{pm}, \ac{cdt}, OTA, \ac{dote} (covered programs) \\
  \midrule
  \ac{emd} & Full-up system \ac{dte}, \ac{lfte} planning, \ac{oa}, cyber \ac{dte} and \ac{ote} integration & Design stability, \ac{trr} success, \ac{oa} results, \ac{lfte} strategy, cyber accreditation & Milestone~C / \ac{lrip} decision & \ac{pm}, \ac{cdt}, OTA, \ac{dote}, \ac{ldto}, warfare centers \\
  \midrule
  Milestone~C / \ac{lrip} & \ac{iote}, \ac{lfte}, cyber \ac{ote}, interoperability certification & \ac{iote} adequacy, \ac{lfte} survivability/lethality, cyber resilience, JITC certification & \ac{frpdr} eligibility & \ac{mdauthority}, \ac{dote}, OTA, \ac{comoptevfor}, JITC, \ac{pm} \\
  \midrule
  \ac{iote} / \ac{lfte} & Independent operational evaluation, realistic threat testing, cyber adversarial assessment & Operational effectiveness/suitability, survivability/lethality, cyber mission assurance & Beyond-\ac{lrip} / \ac{frp} & \ac{dote}, OTA, \ac{comoptevfor}, \ac{pm}, \ac{mdauthority} \\
  \midrule
  \ac{frp} / Full Deployment & \ac{fote}, software increment testing, cyber update validation, sustainment verification & Corrective action validation, \ac{cpd} compliance, cyber resilience, supportability & Full deployment & \ac{mdauthority}, OTA, \ac{pm}, \ac{cdt} \\
  \midrule
  Post-fielding / \ac{fote} & \ac{fote}, software increments, cyber updates, reliability growth & Ongoing effectiveness, deficiency correction, cyber posture, fleet feedback & Sustainment decisions & \ac{pm}, OTA, \ac{cdt}, fleet users \\
  \bottomrule
\end{longtblr}

\subsection{Test Type Matrix}
Different test types answer distinct questions across the \ac{mca} lifecycle~\autocite{DoDI5000-89,DoDI5000-98,DoDM5000-103}.

\begin{longtblr}[
    label   = {tab:test-type-matrix},
    caption = {Test Type Matrix for MCA Programs},
    entry   = {Test Type Matrix for MCA Programs},
    remark{\tSource} = {\tabCite{DoDI5000-89,DoDI5000-98,DoDM5000-103,DoDM5000-96}}
  ]
  {colspec = {@{} X[1.0] X[1.8] X[1.0] X[1.4] X[1.1] X[1.0] @{}}}
  \toprule
  {Test Type} & {Primary Question Answered} & {Typical Timing} & {Requirements / Planning Considerations} & {Major Players} & {MCA Decision Relevance} \\
  \midrule
  Contractor Test \& Evaluation (CT\&E) & Does the design meet contract specifications? & Continuous through \ac{emd} & Contract requirements, quality assurance, data rights & Contractor, \ac{pm} & Design maturity, risk reduction \\
  \midrule
  Developmental Test \& Evaluation (DT\&E) & Did we build the system right against technical requirements? & \ac{msa} through \ac{pd.phase} & \ac{temp} Part~II, \ac{sep}, \ac{ctp}/\ac{tpm} traceability & \ac{pm}, \ac{cdt}, \ac{ldto}, warfare centers & \ac{trr}, \ac{techeval}, \ac{otrr} readiness \\
  \midrule
  Integrated Testing & Can combined DT/OT data support multiple objectives efficiently? & \ac{emd} and \ac{pd.phase} & \ac{temp} integrated test plan, data sharing agreements, independence safeguards & \ac{cdt}, OTA, \ac{pm} & Schedule efficiency, data alignment \\
  \midrule
  Early Operational Assessment (EOA) & Is the operational concept viable for requirements refinement? & \ac{msa} / early \ac{tmrr} & \ac{eoa} plan, user involvement, representative scenarios & OTA, \ac{pm}, users & \ac{cdd} refinement, acquisition strategy \\
  \midrule
  Operational Assessment (OA) & Is the system operationally effective enough to proceed? & Late \ac{emd} / pre-\ac{lrip} & \ac{oa} plan, threat-representative environment, logistics support & OTA, \ac{pm}, fleet & Milestone~C / \ac{lrip} input \\
  \midrule
  Operational Test \& Evaluation (OT\&E) & Can warfighters accomplish missions with the system under realistic conditions? & \ac{pd.phase} / \ac{os.phase} & \ac{temp} OT annex, \ac{dote} approval (covered), operational realism & OTA, \ac{dote}, \ac{comoptevfor}, \ac{pm} & \ac{frp} decision, Title~10 adequacy \\
  \midrule
  Initial Operational Test \& Evaluation (IOT\&E) & Is the production-representative system ready for full deployment? & \ac{lrip} to \ac{frp} & \ac{iote} plan, \ac{dote} oversight, statistical confidence & OTA, \ac{dote}, \ac{comoptevfor} & Beyond-\ac{lrip} / \ac{frp} statutory gate \\
  \midrule
  Follow-on Operational Test \& Evaluation (FOT\&E) & Do corrective actions, software drops, or updates maintain operational effectiveness? & Post-\ac{frp} / major updates & \ac{fote} plan, deficiency validation, cyber update testing & OTA, \ac{pm}, fleet & Sustainment decisions, fielding authorization \\
  \midrule
  Live Fire Test \& Evaluation (LFT\&E) & Is the system survivable and lethal under realistic combat threats? & \ac{emd} through \ac{frp} & \ac{lfte} strategy, Title~10 waiver process, threat surrogates & \ac{dote}, \ac{pm}, test ranges & \ac{frp} statutory requirement (covered systems) \\
  \midrule
  Cybersecurity DT\&E & Is the system cyber-resilient against mission-focused threats throughout development? & Continuous \ac{dte} phases & DoDM~5000.103, \ac{cvpa}, \ac{rmf}, cyber \ac{temp} annex & \ac{pm}, cyber certifiers, \ac{cdt} & \ac{trr}, \ac{otrr}, ATO/CATO \\
  \midrule
  Cybersecurity OT\&E & Can the system maintain mission assurance under cyber attack in operational conditions? & \ac{iote} and \ac{fote} & Cyber \ac{ote} plan, adversarial assessment, mission thread testing & OTA, cyber certifiers, \ac{pm} & \ac{frp}, fielding, cyber accreditation \\
  \midrule
  Interoperability T\&E / JITC Certification & Can the system exchange and use information with joint/coalition/external systems? & \ac{emd} through \ac{frp} & Interface Control Documents, JITC test plan, certification timeline & JITC, \ac{pm}, interface owners & \ac{frp}, fielding, coalition operations \\
  \midrule
  Modeling \& Simulation / VV\&A & Can accredited models supplement or extend live-test evidence? & Lifecycle-wide, especially early and for \ac{lfte} & DoDM~5000.102, \ac{vva} plan, pedigree documentation & \ac{pm}, \ac{cdt}, \ac{dote}, model sponsors & \ac{lfte} waivers, \ac{ote} augmentation, cost savings \\
  \bottomrule
\end{longtblr}

\subsection{ACAT-Level Differences}
\ac{acat} level affects oversight intensity, \ac{mdauthority} level, required documentation, approval authority, and reporting pathways~\autocite{DoDI5000-85}. Non-\ac{dote} oversight does not eliminate the requirement for adequate tailored \ac{te}; statutory, cyber, interoperability, and \ac{lfte} requirements apply regardless of \ac{acat}~\autocite{USC-10-2399,USC-10-2366,DoDI5000-90}.

\begin{longtblr}[
    label   = {tab:acat-te-differences},
    caption = {ACAT-Level Differences for Test and Evaluation},
    entry   = {ACAT-Level Differences for Test and Evaluation},
    remark{\tSource} = {\tabCite{DoDI5000-85,DoDI5000-89,DoDI5000-98}}
  ]
  {colspec = {@{} X[0.7] X[1.4] X[1.2] X[1.3] X[1.2] X[1.3] @{}}}
  \toprule
  {ACAT Level} & {Oversight Intensity} & {MDA Level} & {Required Documentation} & {Approval Authority} & {Reporting} \\
  \midrule
  ACAT~ID & Highest; \ac{dote} statutory oversight & \ac{dae} (USW(A\&S)) & Full \ac{temp}, \ac{dote} reports, SAR, UCR & \ac{dote} approves \ac{temp} and OT plans & Quarterly DAES, SAR, UCR, \ac{dote} reports to Congress \\
  \midrule
  ACAT~IC & High; \ac{dote} oversight for covered programs & \ac{cae} (ASN(RD\&A)) & Full \ac{temp}, \ac{dote} reports (if covered), SAR, UCR & \ac{cae} approves \ac{temp}; \ac{dote} approves OT (covered) & Quarterly DAES, SAR, UCR, \ac{dote} reports (if covered) \\
  \midrule
  ACAT~IB & Moderate; tailored \ac{dote} oversight & \ac{cae} or designee & Tailored \ac{temp}, DAES & \ac{cae} or designee approves \ac{temp} & DAES (if designated) \\
  \midrule
  ACAT~II & Moderate; Service/component oversight & \ac{cae} or PEO & Tailored \ac{temp} & Component Acquisition Executive or PEO & Component reporting \\
  \midrule
  ACAT~III / IV & Low; tailored oversight & \ac{cae}, PEO, or SYSCOM & Tailored \ac{temp} or T\&E Strategy & Designated \ac{mdauthority} & Minimal reporting \\
  \midrule
  Non-DOT&E oversight & Tailored; Service/component channels & Component \ac{mdauthority} & Tailored \ac{temp} or T\&E Strategy & Component Acquisition Executive or designee & Component reporting; statutory requirements still apply \\
  \bottomrule
\end{longtblr}

\note{Board trap: Non-\ac{dote} oversight does not mean no formal \ac{te}. It means \ac{te} is tailored and approved through Service/component channels unless statutory, cyber, interoperability, \ac{lfte}, or other oversight requirements apply~\autocite{DoDI5000-85,DoDI5000-90}.}

\subsection{Major Players and Responsibilities}
\ac{te} success depends on clear role delineation and collaboration across the acquisition enterprise~\autocite{DoDI5000-89,DoDI5000-98}.

\begin{longtblr}[
    label   = {tab:te-players},
    caption = {Major Test and Evaluation Players and Responsibilities},
    entry   = {Major Test and Evaluation Players and Responsibilities},
    remark{\tSource} = {\tabCite{DoDI5000-89,DoDI5000-98,DoDM5000-103,OPTEVFOR-official}}
  ]
  {colspec = {@{} X[1.0] X[1.6] X[1.2] X[1.5] @{}}}
  \toprule
  {Player} & {Role} & {Controls / Approves} & {Does Not Control} \\
  \midrule
  Program Manager (PM) & Overall program responsibility; integrates \ac{te} into acquisition strategy & \ac{temp}, test resources, contractor test plans, deficiency resolution & Independent OT judgments, \ac{dote} reports, statutory certification decisions \\
  \midrule
  Milestone Decision Authority (MDA) & Approves phase entry, \ac{temp}, and acquisition strategy & Milestone decisions, \ac{apb}, \ac{temp} approval, \ac{acat} designation & Test execution details, independent OT findings, technical baseline content \\
  \midrule
  Chief Developmental Tester (CDT) & Orchestrates integrated \ac{dte} schedule, data requirements, and deficiency resolution & Integrated test plan, \ac{dte} resource allocation, deficiency tracking & OT plans, OT execution, \ac{dote} oversight, certification authority \\
  \midrule
  T\&E WIPT / Integrated Test Team & Coordinates \ac{te} planning and execution across stakeholders & Test schedules, data sharing, cross-discipline coordination & Milestone decisions, budget authority, independent OT independence \\
  \midrule
  DOT&E & Independent statutory oversight for covered programs; reports to \ac{secdef} and Congress & OT plan approval, \ac{lfte} strategy approval, OT adequacy assessment & DT planning, \ac{dte} execution, certification decisions, program acquisition strategy \\
  \midrule
  USD(R\&E) / DTE\&A & Sets \ac{te} policy, designates \acp{cdt}, charters \acp{ldto} & \ac{te} policy, \ac{cdt} designations, \ac{ldto} charters, test infrastructure & Individual program \ac{temp} content, OT execution, milestone decisions \\
  \midrule
  Operational Test Agency (OTA) & Plans and executes independent \ac{ote} & OT plan, OT execution, OT report, \ac{otrr} certification & DT planning, contractor testing, program acquisition strategy, certification decisions \\
  \midrule
  OPTEVFOR (Navy) & Navy independent operational test agency for maritime systems & Navy OT plans, OT execution, OT reports, \ac{otrr} certification & DT planning, SYSCOM engineering decisions, certification decisions beyond OT \\
  \midrule
  JITC & Joint interoperability certification for information systems & Interoperability test plans, certification decisions & OT execution, DT planning, program acquisition strategy \\
  \midrule
  Requirements Sponsor / User Community / Fleet & Defines operational context, validates mission effectiveness & \ac{icd}/\ac{cdd}, \ac{conops}, operational scenarios, user feedback & Test planning, technical specifications, certification details \\
  \midrule
  SYSCOM Developmental Test Organizations & Executes \ac{dte}, provides labs/ranges/test infrastructure & DT execution, range scheduling, instrumentation, certification testing & OT planning, OT execution, acquisition strategy \\
  \midrule
  Contractor & Executes CT\&E, provides test articles, supports \ac{dte} & Contractor test plans, test article delivery, CT\&E data & OT plans, OT execution, independent OT judgments, \ac{dote} oversight \\
  \midrule
  Threat / Intelligence Community & Provides threat data, intelligence support, and threat-representative targets & Threat assessments, threat surrogates, intelligence products & Test planning, test execution, certification decisions \\
  \midrule
  Certification Authorities (Cyber, Interoperability, Safety, Airworthiness, Spectrum, Nuclear Surety) & Issue statutory or regulatory certifications & Certification criteria, certification decisions, ATO/CATO & Test planning, test execution, acquisition strategy \\
  \bottomrule
\end{longtblr}

\subsection{Navy-Specific Test and Evaluation Context}
Navy \ac{te} execution reflects maritime operational requirements, shipbuilding constraints, and the distributed nature of naval forces~\autocite{edo-3-5-6-test-eval-part-ii-2025,OPTEVFOR-official}.

\begin{description}
  \item[\textbf{\ac{comoptevfor} as Navy OTA.}] \ac{comoptevfor} is the Navy independent operational test agency, responsible for planning and executing \ac{ote} for maritime systems~\autocite{OPTEVFOR-official}. \ac{comoptevfor} provides independent operational effectiveness and suitability assessments to support \ac{frp} decisions and reports to \ac{dote} for covered programs. The \ac{comoptevfor} commander retains authority to recommend restricted Fleet release if Category~I deficiencies remain~\autocite{edo-3-5-6-test-eval-part-ii-2025}.
  \item[\textbf{SYSCOM Support for DT.}] NAVSEA, NAVWAR, and NAVAIR SYSCOMs support \ac{dte} through warfare centers, engineering sites, labs, ranges, and test infrastructure~\autocite{edo-3-5-5-test-eval-part-i-2025}. SYSCOMs provide technical expertise, instrumentation, certification testing (airworthiness, safety, information assurance, weapon certification), and integration support. SYSCOM developmental test organizations execute \ac{dte} under \ac{cdt} and \ac{ldto} direction but do not replace independent OT.
  \item[\textbf{OPNAV / Resource Sponsors.}] OPNAV N9 warfare-systems sponsors help define requirements, operational context, and portfolio tradeoffs but do not replace independent OT~\autocite{edo-3-5-6-test-eval-part-ii-2025}. N94 is the Test \& Evaluation office within the N9 warfare-systems sponsor family, but N94 is not the sole owner of all Navy testing. N94 coordinates \ac{te} policy and resource alignment across the Navy \ac{te} enterprise.
  \item[\textbf{Interoperability Coordination.}] For systems requiring interoperability with joint or coalition partners, programs must coordinate early with JITC and interface owners to ensure certification timelines align with \ac{frp} decisions~\autocite{JITC-interoperability}. Navy-specific interface control working groups (ICWGs) support this coordination.
  \item[\textbf{Cyber Testing Integration.}] For cyber-heavy systems, cyber \ac{dte} and cyber \ac{ote} must be integrated early in the lifecycle, not deferred until \ac{iote}~\autocite{DoDM5000-103,DoDI5000-90}. Navy cyber testing leverages NIWC Atlantic and Pacific, FCC, and NCDOC for cyber DT\&E and OT\&E support.
\end{description}

\subsection{Common Board Traps}
Avoid these high-yield misconceptions when discussing \ac{mca} testing~\autocite{DoDI5000-85,DoDI5000-89,DoDI5000-98,DoDM5000-103}.

\begin{longtblr}[
    label   = {tab:te-traps},
    caption = {Common Test and Evaluation Board Traps},
    entry   = {Common Test and Evaluation Board Traps},
    remark{\tSource} = {\tabCite{DoDI5000-85,DoDI5000-89,DoDI5000-98,DoDM5000-103}}
  ]
  {colspec = {@{} X[1.8] X[2.2] @{}}}
  \toprule
  {Trap} & {Correct Understanding} \\
  \midrule
  DT\&E proves technical performance; OT\&E proves operational mission performance & \ac{dte} answers ``did we build the system right'' against technical requirements; \ac{ote} answers ``did we build the right capability'' by validating mission effectiveness and suitability under realistic conditions~\autocite{DoDI5000-89,DoDI5000-98} \\
  \midrule
  Integrated testing erases IOT\&E & Combined testing can collect shared data and improve efficiency, but it does not eliminate the need for independent \ac{iote} with operational realism, adequate sample size, and statistical confidence before \ac{frp}~\autocite{DoDI5000-98} \\
  \midrule
  Contractor testing substitutes for independent operational testing & Contractor Test \& Evaluation (CT\&E) is essential for design maturity but is not independent OT; OTA must conduct \ac{iote} with operational users and threat-representative scenarios~\autocite{DoDI5000-89} \\
  \midrule
  ACAT level changes the need for adequate evidence & \ac{acat} level changes oversight intensity and approval pathways, but all programs require adequate \ac{te} evidence for milestone decisions; statutory, cyber, interoperability, and \ac{lfte} requirements apply regardless of \ac{acat}~\autocite{DoDI5000-85,DoDI5000-90} \\
  \midrule
  Cyber testing is a late checklist & Cyber \ac{dte} and cyber \ac{ote} must be integrated throughout the lifecycle, starting in \ac{msa} and continuing through \ac{emd}, \ac{iote}, and \ac{fote}; DoDM~5000.103 mandates lifecycle-wide cyber testing~\autocite{DoDM5000-103,DoDI5000-90} \\
  \midrule
  Interoperability testing is just paperwork certification & Interoperability must be demonstrated in representative operational context through JITC testing, not just documented through interface specifications; certification requires operational validation~\autocite{JITC-interoperability} \\
  \midrule
  LFT\&E applies only to specific systems & Title~10 requires \ac{lfte} for covered systems (crew-occupied platforms, combat vehicles, major munitions, missile defense components) before beyond-\ac{lrip} decisions; waivers require \ac{secdef} approval and congressional notification~\autocite{USC-10-2366,DoDI5000-98} \\
  \midrule
  Modeling and simulation requires no validation & Models used for decision-quality evaluations require appropriate \ac{vva} per DoDM~5000.102; unvalidated models cannot substitute for required live-test evidence without appropriate accreditation~\autocite{DoDM5000-102} \\
  \bottomrule
\end{longtblr}

\subsection{Navy Overlays and Integration}
\ac{opnavinst}~3960-series policy (captured in the coursebook) requires Navy programs to plan \ac{techeval} (led by the engineering agent) before \ac{opeval}, to obtain \ac{opnav} \ac{n9} concurrence on test objectives that bear on Fleet introduction risk, and to route \ac{otrr} packages through \ac{comoptevfor} for certification that logistics, training, and safety enablers are in place~\autocite{edo-3-5-6-test-eval-part-ii-2025}. First-of-class ships also execute Builder's and Acceptance Trials, \ac{insurv} events, and dedicated combat system Ship Qualification Trials that integrate with \ac{techeval}/\ac{opeval} data packages. Navy \ac{cotf} retains authority to recommend restricted Fleet release if Category~I deficiencies remain; those actions feed the Key Roles appendix update for \ac{dasdte}, \ac{cdt}, and \ac{cotf} responsibilities.

\subsection{Fiscal and PPBE Touchpoints}
\ac{rdte} appropriations fund most \ac{dte} events, test articles, range time, and contractor support, while procurement colors such as \ac{scn} or Weapons Procurement, Navy buy dedicated \ac{optest} assets when production-representative hardware is required for statutory adequacy~\autocite{DoDFMR-Vol3,edo-3-5-5-test-eval-part-i-2025}. Operations and Maintenance funds cover Fleet participation, temporary duty, and sustainment of government test ranges. Because \ac{optest} adequacy is a precondition for \ac{frp} and sustainment decisions, Program Objective Memorandum issue papers must protect test range modernization, threat-representative targets, and Modeling \& Simulation accreditation costs or risk a \ac{ppbe} disconnect when \ac{comoptevfor} identifies unfunded test objectives.

\subsection{Assumptions and Uncertainties}
\begin{itemize}
  \item Coursebook Modules~3.5.5 and~3.5.6 (March~2025 editions) remain the latest \ac{don}-specific \ac{te} references.
  \item \ac{comoptevfor} manual updates post-2024 were not available offline; Navy-specific terminology (e.g., \ac{opeval} taxonomy) is based on coursebook guidance.
  \item Upcoming \ac{ndaa} changes to Title~10 \ac{te} statutes have not been incorporated pending official publication.
\end{itemize}

%====================
% * End of file
%====================
\IfSubfilesClassLoaded{
  \RenewDocumentCommand{\entryname}{}{\textbf{\color{Modern} Acronym}}
  \RenewDocumentCommand{\descriptionname}{}{\textbf{\color{Modern} Definition}}
  \printnoidxglossary[
    type=\acronymtype,
    title=Acronyms,
    style=long-booktabs
  ]}{}
\end{document}
